\documentclass[11pt]{article}
\usepackage[margin=1in]{geometry}
\usepackage{amsmath,amssymb,amsthm}
\usepackage{graphicx}
\usepackage{hyperref}
\usepackage{microtype}
\hypersetup{hidelinks}

\newtheorem{theorem}{Theorem}
\newtheorem{proposition}[theorem]{Proposition}
\newtheorem{corollary}[theorem]{Corollary}
\newtheorem{remark}[theorem]{Remark}
\newcommand{\R}{\mathbb R}
\newcommand{\N}{\mathbb N}
\newcommand{\K}{\mathcal K}
\newcommand{\Hh}{\mathcal H}

\title{A Canonical Lower-Semicontinuous Kuratowski Extension\\
and the Sharp Obstruction to Affine Interpolation}
\author{Candidate solution packet for arXiv:2402.11284}
\date{11 August 2026}

\begin{document}
\maketitle

\begin{center}
\fbox{\parbox{0.91\textwidth}{\textbf{Status.}
Candidate full two-sided resolution; likely valid, pending human review.
Every real-valued lower-semicontinuous datum has an explicit finite extension;
for nonnegative data there is a maximal chord-subaffine extension.  The exact
endpoint-affine analogue used in the Lipschitz-free-space proof is impossible
in general.}}
\end{center}

\section{Source question}

The source is D.~Bate, S.~Eriksson-Bique, and E.~Soultanis,
\emph{Fragment-wise differentiable structures}, arXiv:2402.11284v1
\cite{BES}.  In Remark~2.1 on PDF page~7, the authors explain that their
natural convex extension is not well-defined after the standard Kuratowski
embedding and ask whether a canonical extension preserving semicontinuity
exists there.

\begin{figure}[h]
\centering
\includegraphics[width=\textwidth]{figures/open_problem_crop.png}
\caption{Remark~2.1 in the source paper (PDF page~7).}
\end{figure}

The word ``canonical'' is not axiomatized in the remark.  We resolve the
question using the strongest natural one-sided interpolation property and
then show that the corresponding equality property is impossible.

\section{Definitions and statement}

Let $(X,d)$ be a separable metric space, fix $o\in X$, and let
$D\subset X$ be dense.  The base-pointed Kuratowski embedding is
\[
  \K:X\longrightarrow \ell^\infty(D),\qquad
  \K(x)(s)=d(x,s)-d(o,s).
\]
It is an isometry.  Set
\[
 \Hh_{\K}(X)=\bigcup_{x,y\in X}[\K(x),\K(y)].
\]
All topology below is the relative norm topology on this semihull.

For any $\rho:X\to[0,\infty]$, define its raw two-point roof by
\begin{equation}\label{eq:raw}
 R_\rho(z)=\inf\bigl\{(1-t)\rho(a)+t\rho(b):
 z=(1-t)\K(a)+t\K(b),\ a,b\in X,\ t\in[0,1]\bigr\}.
\end{equation}
We use the standard convention $0\cdot\infty=0$.  Define the relaxed roof
\begin{equation}\label{eq:relaxed}
 E_\rho(z)=\sup_{r>0}\ \inf\bigl\{R_\rho(w):
 w\in\Hh_{\K}(X),\ \|w-z\|_\infty<r\bigr\}.
\end{equation}

\begin{theorem}[Canonical Kuratowski extension]\label{thm:main}
If $\rho:X\to[0,\infty]$ is lower semicontinuous, then $E_\rho$ is a
lower-semicontinuous extension of $\rho$:
\[
 E_\rho(\K(x))=\rho(x)\qquad(x\in X).
\]
Moreover:
\begin{enumerate}
\item for all $a,b\in X$ and $t\in[0,1]$,
\begin{equation}\label{eq:subaffine}
E_\rho((1-t)\K(a)+t\K(b))
\le (1-t)\rho(a)+t\rho(b);
\end{equation}
\item $E_\rho$ is the greatest lower-semicontinuous function
$g:\Hh_{\K}(X)\to[0,\infty]$ satisfying every inequality
\eqref{eq:subaffine};
\item the assignment $\rho\mapsto E_\rho$ is order preserving, positively
homogeneous, and preserves nonnegative constant functions.
\end{enumerate}
Thus \eqref{eq:relaxed} is canonical by a maximal universal property.
\end{theorem}

\begin{corollary}[Literal real-valued extension]\label{cor:real}
Assume additionally that $X$ is complete.  Every lower-semicontinuous
$\rho:X\to\R$ admits an explicit finite lower-semicontinuous extension
$T_\rho:\Hh_{\K}(X)\to\R$ with
$T_\rho\circ\K=\rho$.
\end{corollary}

\begin{theorem}[Sharp obstruction to the free-space formula]\label{thm:c4}
There is a four-point metric space $X$ and a continuous
$\rho:X\to[0,1]$ for which no function
$G:\Hh_{\K}(X)\to[0,\infty]$ can satisfy
\[
G((1-t)\K(a)+t\K(b))=(1-t)\rho(a)+t\rho(b)
\]
simultaneously for every $a,b\in X$ and $t\in[0,1]$.
Consequently no general Kuratowski extension scheme can retain the exact
endpoint-affine identity of the Lipschitz-free-space construction.
\end{theorem}

\section{Proof intuition}

The raw roof removes ambiguity by choosing the cheapest value among all chord
representations.  Its lower-semicontinuous relaxation is automatically the
largest lower-semicontinuous function compatible with all resulting upper
bounds.  Relaxation could in principle lower the boundary data.  The special
geometry of the Kuratowski embedding prevents that: evaluating a barycenter
near $\K(x)$ at the coordinate $x$ shows that almost all of its barycentric
mass must lie on endpoints near $x$.  Lower semicontinuity of $\rho$ then
forces the cost back up to $\rho(x)$.

The obstruction is the cycle metric on four points.  Opposite chords have the
same midpoint, so exact affine interpolation assigns two incompatible values
there.  This explains precisely why the free-space construction can be
replaced by a canonical inequality but not by the exact identity used for
weighted gap integrals.

For arbitrary real-valued data, first compress the values monotonically into
the interval $(1,2)$, use the roof theorem there, and keep the extension away
from the endpoints of that interval by a canonical distance-to-$\K(X)$ floor.
Completeness makes $\K(X)$ closed, so the floor is strictly above $1$ off the
embedded space.  Inverting the compression yields a finite real-valued
lower-semicontinuous extension.

\section{Proofs}

We first isolate the Kuratowski estimate that makes exact boundary recovery
possible.

\begin{proposition}\label{prop:exposure}
For $a,b,x\in X$ and $t\in[0,1]$,
\begin{equation}\label{eq:exposure}
\bigl\|(1-t)\K(a)+t\K(b)-\K(x)\bigr\|_\infty
\ge (1-t)d(a,x)+t d(b,x).
\end{equation}
\end{proposition}

\begin{proof}
If $x\in D$, evaluate the vector on the left at $s=x$.  The base-point terms
cancel, and the value is the nonnegative number on the right.  For general
$x$, choose $s_n\in D$ with $s_n\to x$.  The coordinate values converge to
the same number, so their supremum is at least that limit.
\end{proof}

\begin{proof}[Proof of Theorem~\ref{thm:main}]
For any extended-real function on a metric space, the expression in
\eqref{eq:relaxed} is its lower-semicontinuous relaxation, hence $E_\rho$ is
lower semicontinuous and $E_\rho\le R_\rho$.  Formula \eqref{eq:raw} then
immediately gives \eqref{eq:subaffine}.

It remains to prove the boundary trace.  Fix $x\in X$.  The representation
$\K(x)=(1-t)\K(x)+t\K(x)$ shows
$E_\rho(\K(x))\le R_\rho(\K(x))\le\rho(x)$.

Suppose first that $0<\rho(x)<\infty$, and fix $c<\rho(x)$ with $c\ge0$.
By lower semicontinuity there is $\delta>0$ such that
$\rho(u)>c$ whenever $d(u,x)<\delta$.  Let
$w=(1-t)\K(a)+t\K(b)$ and put $r=\|w-\K(x)\|_\infty$.
By Proposition~\ref{prop:exposure}, the total barycentric weight of endpoints
outside $B(x,\delta)$ is at most $r/\delta$.  Nonnegativity therefore gives
\[
 (1-t)\rho(a)+t\rho(b)\ge c(1-r/\delta).
\]
This holds for every representation of $w$, so
$R_\rho(w)\ge c(1-r/\delta)$.  Taking the infimum over the radius-$r_0$
neighborhood of $\K(x)$ and then letting $r_0\downarrow0$ yields
$E_\rho(\K(x))\ge c$.  Letting $c\uparrow\rho(x)$ proves equality.  If
$\rho(x)=0$, the lower bound is automatic.  If $\rho(x)=\infty$, repeat the
argument with an arbitrary finite $c$; then let $c\to\infty$.

For maximality, let $g$ be lower semicontinuous and satisfy all the chord
inequalities.  Taking the infimum over representations gives $g\le R_\rho$.
For every $z$,
\[
g(z)=\sup_{r>0}\inf_{\|w-z\|_\infty<r}g(w)
\le \sup_{r>0}\inf_{\|w-z\|_\infty<r}R_\rho(w)=E_\rho(z),
\]
where all points $w$ lie in $\Hh_{\K}(X)$.  This proves the universal
property.  Monotonicity, positive homogeneity, and preservation of constants
follow directly from \eqref{eq:raw} and \eqref{eq:relaxed}.
\end{proof}

\begin{proof}[Proof of Corollary~\ref{cor:real}]
Fix the increasing homeomorphism
\[
h:\R\longrightarrow(1,2),\qquad
h(s)=\frac32+\frac1\pi\arctan s.
\]
Apply Theorem~\ref{thm:main} to the lower-semicontinuous function
$f=h\circ\rho$.  Since every endpoint value lies strictly between $1$ and
$2$, its relaxed roof satisfies
\[
1\le E_f(z)<2\qquad(z\in\Hh_{\K}(X)).
\]
The image $\K(X)$ is closed because $X$ is complete and $\K$ is an isometry.
Put
\[
q(z)=1+\frac{\operatorname{dist}(z,\K(X))}
                 {1+\operatorname{dist}(z,\K(X))},
\qquad
T_\rho(z)=h^{-1}\!\left(\max\{E_f(z),q(z)\}\right).
\]
The maximum is lower semicontinuous, is strictly between $1$ and $2$ at every
point, and equals $f(x)$ at $z=\K(x)$: there $q=1$ and
$E_f(\K(x))=f(x)>1$.  Since $h^{-1}$ is continuous and increasing,
$T_\rho$ is finite and lower semicontinuous, and
$T_\rho(\K(x))=\rho(x)$.
\end{proof}

\begin{proof}[Proof of Theorem~\ref{thm:c4}]
Let $X=\mathbb Z/4\mathbb Z$ with the cycle metric
$d(i,j)=\min\{|i-j|,4-|i-j|\}$, and use base point $0$ and all four points as
coordinates.  The distance rows are
\[
 D_0=(0,1,2,1),\quad D_1=(1,0,1,2),\quad
 D_2=(2,1,0,1),\quad D_3=(1,2,1,0).
\]
Since $\K(i)=D_i-D_0$, the identity $D_0+D_2=D_1+D_3$ gives
\[
 \frac{\K(0)+\K(2)}2=\frac{\K(1)+\K(3)}2=:m.
\]
Set $\rho(0)=\rho(2)=0$ and $\rho(1)=\rho(3)=1$.  Every function on this
finite metric space is continuous.  Exact affine interpolation on the first
opposite chord forces $G(m)=0$, while interpolation on the second forces
$G(m)=1$, a contradiction.
\end{proof}

\section{Interpretation and limitations}

Corollary~\ref{cor:real} answers the literal existence question for all
real-valued lower-semicontinuous data under the completeness hypothesis of the
source paper.  Theorem~\ref{thm:main} gives the stronger order-theoretic answer
for the nonnegative integrands used in the application.  It is nontrivial on
the semihull (unlike the extension that assigns $+\infty$ off $\K(X)$), and it
is forced as the greatest lower-semicontinuous chord-subaffine extension.

It does not reproduce the identity
\[
 \bar\rho((1-t)\K(a)+t\K(b))=(1-t)\rho(a)+t\rho(b)
\]
along every gap chord.  Theorem~\ref{thm:c4} shows that this limitation cannot
be removed for arbitrary metric spaces.  Therefore the positive construction
does not by itself replace the exact line-integral identity in Section~3.1 of
the source paper.

\section{Verification and novelty bounds}

The proof is self-contained.  The finite script in the packet independently
checks the cycle distance matrix, the midpoint identity, and the conflicting
endpoint averages using exact integer arithmetic.  It is a check of the
example, not a substitute for the general proof.

On 11 August 2026, exact arXiv-id/title/core-keyword searches of this run's
registry and result indexes found no prior result.  Exact-phrase and nearby
keyword searches of the locally parsed arXiv corpus found the question only
in arXiv:2402.11284.  Locally available citing sources arXiv:2504.16657,
2508.08017, and 2510.25715 do not discuss the Kuratowski-extension issue.
Bounded external web searches returned no usable records in this environment;
novelty confidence is therefore moderate.  No external theorem is used.

\textbf{Human-review recommendation.}  Check the extended-value boundary
argument and decide whether the maximal chord-subaffine universal property is
the intended formalization of ``canonical'' in Remark~2.1.  Independently,
the four-cycle impossibility theorem is unconditional and exact.

\begin{thebibliography}{9}
\bibitem{BES}
D.~Bate, S.~Eriksson-Bique, and E.~Soultanis,
\emph{Fragment-wise differentiable structures}, arXiv:2402.11284v1 (2024).
\end{thebibliography}

\end{document}
